\documentclass[11pt]{article}

\input{preambles/header}
\input{preambles/notation}

% this sets the paths for graphicsx to search for figures
\graphicspath{{../Empirics/Output/}{./}}
%\renewcommand{\ignore}[1]{ }

\setcounter{tocdepth}{3}
\setcounter{secnumdepth}{3}
\clubpenalty = 10000
\widowpenalty = 10000


\setlength{\parskip}{0.2mm}

\usepackage{pgfplotstable}

\usepackage{geometry}
\geometry{left=1in}
\geometry{right=1 in}
\geometry{top=1in}
\geometry{bottom=1in}
\linespread{1.4}


\usepackage{caption}
\captionsetup{labelformat=empty,labelsep=none}

\externaldocument{online_appendix}

\begin{document}
\thispagestyle{empty}


% \vspace{20pt}
\begin{center}
% \textcolor{white}{fill}\\
% \vspace{20pt}
\LARGE{GENERATIVE AI AT WORK$^{*}$}%\\\mbox{ }
\end{center}


\begin{singlespace}
\setstretch{1.5}
\begin{center}
\textsc{Erik Brynjolfsson}\\
\textsc{Danielle Li}\\
\textsc{Lindsey Raymond}
\end{center}

\end{singlespace}


% \begin{center}
% \vspace{.20in}
% \normalsize{November 2024}
%\normalsize{\today}


% \vspace{10pt}
% First Draft: 23 April 2023\\
% \end{center}
% \vspace{10pt}

\begin{abstract}
% \vspace{-11pt}
% \singlespacing
% \noindent


We study the staggered introduction of a generative AI-based conversational assistant using data from 5,172 customer support agents.  Access to AI assistance increases worker productivity, as measured by issues resolved per hour, by 15 percent on average, with substantial heterogeneity across workers.  The effects vary significantly across different agents. Less experienced and lower-skilled workers improve both the speed and quality of their output while the most experienced and highest-skilled workers see small gains in speed and small declines in quality.  We also find evidence that AI assistance facilitates worker learning and improves English fluency, particularly among international agents. While AI systems improve with more training data, we find that the gains from AI adoption are largest for moderately rare problems, where human agents have less baseline experience but the system still has adequate training data. Finally, we provide evidence that AI assistance improves the experience of work along several dimensions: customers are more polite and less likely to ask to speak to a manager. \textit{JEL} Codes: D80, J24, M15, M51, O33.
\end{abstract}


% \indent \textbf{Keywords}: Generative AI, Large Language Models, Technology Adoption, Worker Productivity, Worker Learning, Experience of Work, Organizational Design. 

% \bigskip


% \hspace{-25pt}\rule[0.1ex]{0.33\textwidth}{0.2mm}\vspace{0.05in}\\\indent
\let\thefootnote\relax\footnote{$^{*}$%This paper is part of the National Bureau of Economic Research (NBER)'s working paper series, \#WP31161. 
We gratefully acknowledge the coeditors, Lawrence Katz and Andrei Shleifer, and anonymous referees for many valuable insights and suggestions. We are grateful to Daron Acemoglu, David Autor, Amittai Axelrod, Eleanor Dillon, Zayd Enam, Luis Garicano, Alex Frankel, Sam Manning, Sendhil Mullainathan, Emma Pierson, Scott Stern, Ashesh Rambachan, John Van Reenen, Raffaella Sadun, Kathryn Shaw, Christopher Stanton, Sebastian Thrun, and various seminar participants for helpful comments and suggestions. We thank Max Feng and Esther Plotnick for providing excellent research assistance and the Stanford Digital Economy Lab for financial support. The content is solely the responsibility of the authors and does not necessarily represent the official views of Stanford University, MIT, or the NBER. }
% \end{singlespace}

% \thispagestyle{empty}
\linespread{1.5}
\setcounter{page}{0}

% \clearpage
\newpage
\section{Introduction}
The emergence of generative artificial intelligence (AI) has attracted significant attention for its potential economic impact.  Although various generative AI tools have performed well in laboratory settings, questions remain about their effectiveness in real-world settings, where they may encounter unfamiliar problems, face organizational resistance, or provide misleading information \citep{peng2023check, rooseConversationBingChatbot2023}.

We provide early evidence on the impact of generative AI deployed at scale in the workplace. We study the adoption of a generative AI tool that provides conversational support to customer service agents. We find that access to AI assistance increases the productivity of agents by 15 percent, as measured by the number of customer issues they are able to resolve per hour. The gains accrue disproportionately to less-experienced and lower-skill customer support workers indicating that generative AI systems may be capable of capturing and disseminating the behaviors of the most productive agents.

% A note on terminology. There are many definitions of artificial intelligence and of intelligence itself---\citet{legg2007collection} list over 70 of them.  In this paper, we define ``artificial intelligence'' (AI) as an umbrella term that refers to systems that exhibit intelligent behavior, such as learning reasoning and problem-solving.  ``Machine learning'' (ML) is a branch of AI that uses algorithms to learn from data, identify patterns, and make predictions or decisions without being explicitly programmed metha \citep{aivsml}. Large language models (LLMs) and tools built around LLMs such as ChatGPT are an increasingly important application of machine learning. LLMs generate new content, making them a form of ``generative AI.'' ``Generative AI'' is a type of artificial intelligence that can create new content, such as text, images, or music, by learning patterns from existing data.}

Computers and software have transformed the economy with their ability to perform certain tasks with far more precision, speed, and consistency than humans.  To be effective, these systems typically require explicit and detailed instructions for how to transform inputs into outputs: a software engineer must program computers.  Yet, despite significant advances in traditional computing, many workplace activities, such as writing emails, analyzing data, or creating presentations, are difficult to codify and have therefore defied computerization.
%\footnote{Tacit knowledge refers to the knowledge and skills that individuals possess but are unable to express explicitly. It is often intuitive and nonverbal, acquired through personal experiences, observations, and practice over time. Tacit knowledge is deeply ingrained in an individual's behavior and can be difficult to transfer or convey to others through traditional methods such as training or manuals \citep{polanyi_tacit_1966}. }

Machine learning (ML) algorithms work differently from traditional computer programs. Instead of requiring explicit instructions, machine learning algorithms infer instructions from examples. Given a training set of images, for instance, ML systems can learn to recognize specific individuals without requiring a list of the physical features that identify a given person. As a result, ML systems can perform tasks even when no instructions exist \citep{polanyi_tacit_1966, autorPolanyiParadoxShape2014, brynjolfssonmitchell2017}. The data to train these ML systems is often generated by human workers, who naturally vary in their abilities. As a result, ML tools may differentially introduce lower-performing workers to new skills and techniques, causing varied productivity shifts even among workers engaged in the same task.

%The use of ML tools may therefore differentially expose lower-performing workers to new skills and techniques, leading to disparate changes in productivity even among workers performing the same task. 
% \footnote{As \citet{meijer2018behind} puts it, ``where the Software 1.0 Engineer formally specifies their problem, carefully designs algorithms, composes systems out of subsystems or decomposes complex systems into smaller components, the Software 2.0 Engineer amasses training data and simply feeds it into an ML algorithm...''}


%Drawing on many examples of tasks---making sales pitches, driving a truck, or diagnosing a patient---performed at varying levels of proficiency, the systems can implicitly learn what specific behaviors and characteristics distinguish highly effective workers. The use of ML tools may therefore differentially expose lower-performing workers to new skills and techniques, leading to disparate changes in productivity even among workers performing the same task. 

We study the impact of generative AI on productivity and worker experience in the customer service sector, an industry with one of the highest surveyed rates of AI adoption \citep{chui_global_2021}. We examine the staggered deployment of a chat assistant, using data from 5,172 customer support agents working for a Fortune 500 firm that sells business process software. The tool we study is built on Generative Pre-trained Transformer 3 (GPT-3), a member of the Generative Pre-trained Transformer family of large language models developed by OpenAI \citep{GPT4TechnicalReport2023}.  The AI system monitors customer chats and provides agents with real-time suggestions for how to respond. The AI system is designed to augment agents, who remain responsible for the conversation and are free to ignore or edit the AI's suggestions.

We have four sets of findings.  

First, AI assistance increases worker productivity, resulting in a 15 percent increase in the number of chats that an agent successfully resolves per hour. These productivity increases are based on shifts in three components of overall productivity: a decline in the time it takes an agent to handle an individual chat, an increase in the number of chats that an agent handles per hour (agents may handle multiple chats at once), and a small increase in the share of chats that are successfully resolved.  

Second, the impact of AI assistance varies widely among agents.  Less-skilled and less-experienced workers improve significantly across all productivity measures, including a 30 percent increase in the number of issues resolved per hour.  The AI tool helps also newer agents move more quickly down the experience curve: treated agents with two months of tenure perform just as well as untreated agents with more than six months of tenure.  In contrast, AI has little impact on the productivity of higher-skilled or more experienced workers.  Indeed, we find evidence that AI assistance leads to a small \textit{decrease} in the quality of conversations conducted by the most skilled agents.  These results contrast, in spirit, with studies that find evidence of skill-biased technical change for earlier waves of computer technology and robotics \citep{bres2002, bartel2007, dixon2020}. 

Our third set of results investigates the mechanism underlying our main findings. Agents who follow AI recommendations more closely see larger gains in productivity, and adherence rates increase over time.  We also find that experience with AI recommendations can lead to durable learning. Using data on software outages---periods when the AI's output is unexpectedly interrupted due to technical issues---we find that workers see productivity gains relative to their pre-AI baseline even when AI recommendations are unavailable.  The gains are most pronounced among workers who had greater exposure to AI and followed AI suggestions more closely. In addition, we examine the heterogeneous impact of AI access by the technical support conversation topics that agents encounter.  While AI systems improve with access to more training data, we find that the gains to AI adoption---when used to complement human workers---are largest for relatively rare problems, perhaps because agents are already capable of addressing the problems they encounter most frequently. When there is insufficient training data for a topic, the system may not provide suggestions at all.  We further analyze the text of agents' chats and provide evidence that access to AI improves their English language fluency, especially among international agents.  Finally, we compare the text of conversations before and after AI and provide suggestive evidence that AI adoption drives convergence in communication patterns: low-skill agents begin communicating more like high-skill agents.

% .\footnote{The term ``contact center'' updates the term ``call center'' to reflect the fact that a growing proportion of customer service contacts no longer involve phone calls.}
The fourth set of results focuses on the dynamics between agents and customers. Contact center (formerly known as ``call center'') work is often challenging, with agents frequently facing hostile interactions from anonymous, frustrated customers. Many agents also work overnight shifts to align with U.S. business hours due to the prevalence of offshoring. While AI assistance can help agents communicate more effectively, it may risk making agents seem mechanical or inauthentic. Our findings show that access to AI assistance significantly improves customer treatment of agents, as reflected in the tone of customer messages. Customers are also less likely to question agents' competence by asking to speak to a supervisor. Notably, these changes come alongside a decrease in worker attrition, which is driven by the retention of newer workers.  

Our findings show access to generative AI suggestions can increase the productivity of individual workers and improve their experience of work.  We emphasize, however, that these findings capture medium-run impacts within a single firm. Our paper is not designed to shed light on the aggregate employment or wage effects of generative AI tools. In the longer run, firms may respond to increasing productivity among novice workers by hiring more of them or by seeking to develop more powerful AI systems that replace labor altogether.  While the introduction of generative AI may increase demand for lower-skill labor \textit{within} an occupation, the equilibrium response to AI assistance may lead to \textit{across} occupation shifts in labor demand that instead benefit higher-skill workers \citep{autorSkillContentRecent, acemogluLowSkillHighSkillAutomation2018, acemogluSimpleMacroeconomicsAI}. Unfortunately, our data do not allow us to observe changes in wages, overall labor demand, or the skill composition of workers hired at the firm.
% \footnote{For example, transcripts from a fully automated customer service division could be sent to human analysts who mine this data for information on how to improve a product.}  

The results also underscore the longer-term challenges raised by the adoption of AI systems.  In our data, top workers increase their adherence to AI recommendations, even though those recommendations marginally decrease the quality of their conversations.  Yet, with fewer original contributions from the most skilled workers, future iterations of the AI model may be less effective in solving new problems. Our work therefore raises questions about how these dynamics play out over the long run. 

This paper is related to a large literature on the impact of technological adoption on worker productivity and the organization of work \citep[e.g.][]{rosen1981, autorComputingInequalityHave1998, BrynHittJEP2000, atheystern2002, bartel2007, acemoglu2007, bloom2014, michaelsHasIctPolarized2014, garicano2015, hoffman2018, acemogluRobotsJobsEvidence2020, feltenOccupationalHeterogeneityExposure2023}.  Many of these studies, particularly those focused on information technologies, find evidence that IT complements higher-skill or more-educated workers \citep{bres2002,akermanSkillComplementarityBroadband2015,taniguchiICTCapitalSkillComplementarity2022}. For instance, \citet{bartel2007} find that firms that adopt IT tend to use more skilled labor and adoption is associated with increased skill requirements for machine operators in valve manufacturing. Other research investigates how technology affects workers based on their educational attainment and occupation. For example,  \citet{acemogluRobotsJobsEvidence2020} find that the negative effects of robots on employment are most pronounced for workers in blue-collar occupations who lack a college education. 

Fewer studies focus on AI-based technologies, generative or not.  \citet{zolas2020, acemoglu2022,  calvino2023} examine economy-wide data from the US and OECD, finding that the adoption of AI tools is concentrated among larger and younger firms with relatively high productivity. So far, evidence is mixed on the effects of AI on productivity. For example, \citet{acemoglu2022} find no detectable relationship between investments in AI-specific tools and firm outcomes, while \citet{babina2022} find evidence of a positive relationship between firms' AI investments and their subsequent growth and valuations. All the studies stress that determining the productivity effects of AI technologies is difficult because AI-adopting firms differ substantially from non-adopters. 
% \footnote{\citet{oecd2023} reports that when surveyed firms are asked directly, 57\% of employers in finance and 63\% in manufacturing reported that AI positively impacted productivity, and 80\% of surveyed workers who work with AI report higher job performance.} 

Other studies present mixed findings on the effects of AI tools on decision-making, often revealing challenges in human-AI collaboration. As an example of a positive finding, \citet{kanazawa2022} find a non-generative AI tool that suggests customer-rich routes to taxi drivers reduces their search time by 5 percent, with the least experienced drivers benefiting the most. By contrast, several other studies find humans assisted by AI make worse decisions than do either humans or AI alone. \citep{hoffman2018, 
poursabzisangdeh2021, angelova2023, agrawal2023}. In fact, a meta-analysis of more than 100 experimental studies concludes that, on average, human-AI collaborations underperform both the AI alone and the best human decision-makers \citep{vaccaro2024combinationshumansaiuseful}. These results underscore the particular challenges introduced when using AI-based tools designed to augment human decision making.  

In this paper, we provide micro-level evidence on the adoption of a generative AI tool across thousands of workers employed by a given firm and its subcontractors.  Our work is closely related to several other studies examining the impacts of generative AI in lab-like settings.  \citet{peng2023impact} recruit software engineers for a specific coding task (writing an HTTP server in JavaScript) and show those given access to the AI tool GitHub Copilot complete this task twice as quickly as a control group. Similarly, \citet{noyExperimentalEvidenceProductivity2023} conduct an online experiment showing that subjects given access to ChatGPT complete professional writing tasks more quickly. In the legal domain, \citet{choiAIAssistanceLegal2023} show AI assistance helps law students on a law school exam, while in management consulting, \citet{dellacquaNavigatingJaggedTechnological2023} find access to GPT-4 suggestions improves the quality of responses on some management consulting tasks, but can negatively impact performance on tasks outside its capabilities. Consistent with our findings, \citet{noyExperimentalEvidenceProductivity2023}, 
\citet{choiAIAssistanceLegal2023}, \citet{peng2023check} and \citet{dellacquaNavigatingJaggedTechnological2023} find generative AI assistance compresses the productivity distribution, with lower-skill workers driving the bulk of improvements. Our paper, however, provides new evidence of longer-term effects in a real-world workplace where we also track patterns of learning, customer-side effects, and changes in the experience of work.  


\section{Generative AI and Large Language Models}

\subsection{Technical Primer}

The rapid pace of AI development and public release tools such as ChatGPT, GitHub Copilot, and DALL-E have attracted widespread attention, optimism, and alarm \citep{whitehouse2022}.  Such tools are examples of ``generative AI,'' a class of machine learning technologies that can generate new content---such as text, images, music, and video---by analyzing patterns in existing data. This paper focuses on an important class of generative AI, large language models (LLMs). LLMs are neural network models designed to process sequential data \citep{bubeck2023sparks}.  An LLM is trained by learning to predict the next word in a sequence, given what has come before, drawing on a large corpus of text, such as Wikipedia, digitized books, and portions of the Internet. From its knowledge base of the statistical co-occurrence of words, the LLM generates new text that is grammatically correct and semantically meaningful.  Although ``large language model'' implies human language, these techniques can also be used to produce other forms of sequential data (``text'') such as computer code, protein sequences, audio, and chess moves \citep{el2023}. 

Four factors are driving improvements in generative AI: computing scale, earlier innovations in model architecture, the ability to ``pre-train'' using large amounts of unlabeled data, and refinements in training techniques \citep{radfordImprovingLanguageUnderstanding2018, radfordLanguageModelsAre2019,  ouyangTrainingLanguageModels2022, liuSummaryChatGPTGPT42023}.
%\footnote{For a more detailed technical review of progress, see \citet{radfordImprovingLanguageUnderstanding2018, radfordLanguageModelsAre2019, liuSummaryChatGPTGPT42023, ouyangTrainingLanguageModels2022}.}  

The quality of LLMs is strongly dependent on scale: the amount of computing power used for training, the number of model parameters, and the dataset size \citep{kaplan2020scaling}.  The GPT-3 model included 175 billion parameters, was trained on 300 billion tokens, and incurred approximately \$5 million dollars in training costs; the GPT-4 model, meanwhile, includes 1.8 trillion parameters and was trained on 13 trillion tokens at an estimated computing-only cost of \$65 million \citep{liOpenAIGPT3Language2020, brownLanguageModelsAre2020, gpt4scale}.

% transformer models
% \footnote{For instance, a model would keep track of ``the, 1'' instead of only ``the'' (if ``the'' was the first word in the sentence).}  
Modern LLMs use two earlier key innovations in model architecture: positional encoding and self-attention. Positional encodings keep track of the order in which a word occurs in a given input. Self-attention assigns importance weights to each word in the context of the entire input text. Together, these technological advances enable models to capture long-range semantic relationships within an input text, even when that text is broken up into smaller segments and processed in parallel \citep{bahdanau2016neural, vaswaniAttentionAllYou2017}.

These innovations in model architecture enable LLMs to train on large amounts of unlabeled data from sources such as Reddit and Wikipedia. Unlabeled data are far more prevalent than labeled data, allowing LLMs to learn about natural language on a much larger training corpus \citep{brownLanguageModelsAre2020}.  By seeing, for example, that the word ``yellow'' is more likely to be observed with ``banana'' or ``sun'' or ``rubber duckie,'' the model can learn about semantic and grammatical relationships without explicit guidance \citep{radfordImprovingLanguageUnderstanding2018}.  This approach enables LLMs to learn a foundational understanding of language patterns and relationships that can be adapted or further fine-tuned for a specific task.

Fine-tuning can refine general-purpose LLMs output to match the priorities of a specific setting \citep{ouyangTrainingLanguageModels2022, liuSummaryChatGPTGPT42023}. Fine-tuning can help eliminate factually incorrect or inappropriate responses or prioritize a particular tone of response. Such improvements make a general-purpose model better suited to its specific application \citep{ouyangTrainingLanguageModels2022}. For example, a model trained to generate social media content can be further trained on labeled data that contain not just the content of a post or tweet, but also information on the user engagement it attracts. 

Together, these innovations in computing scale, model architecture and training have generated meaningful improvements in model performance. The Generative Pre-trained Transformer (GPT) family of models, in particular, has attracted considerable attention for outperforming humans on tests such as the US legal bar exam \citep{liuSummaryChatGPTGPT42023, bubeck2023sparks, GPT4TechnicalReport2023}.


\subsection{The Economic Impacts of Generative AI}\label{sec:aibackground}

Computers have historically excelled at executing pre-programmed instructions, making them particularly effective at tasks that can be described by explicit rules \citep{autorPolanyiParadoxShape2014}. Consequently, computerization has disproportionately decreased the demand for workers performing routine and repetitive tasks such as data entry, bookkeeping, and assembly line work, reducing wages in these jobs \citep{acemoglu2011skills}. At the same time, computerization has increased the demand for workers who possess complementary skills such as programming, data analysis, and research.  As a result, technology-related shifts in the labor market have contributed to increased wage inequality in the United States and have been linked to a variety of organizational changes \citep{katzChangesRelativeWages1992c, BrynHittJEP2000,  bres2002, 
autorSkillContentRecent, bakerMakeBuyTrucking2003,michaelsHasIctPolarized2014, oecd2023}.

%This type of ``tacit knowledge'' underlies most tasks humans perform, both on and off the job \citep{polanyi_tacit_1966, autorPolanyiParadoxShape2014}.

In contrast, generative AI tools do not require explicit instructions to perform tasks. If asked to write an email denying an employee a raise, generative AI tools will likely respond with a professional and conciliatory note. The model will have seen many examples of workplace communication without the need for a programmer to explicitly define what professional writing looks like. These machine learning methods behind AI enable computers to perform non-routine tasks that rely on tacit knowledge and experience. AI has shown promise on tasks traditionally dominated by highly skilled professionals insulated from prior waves of automation, including complex mathematics, scientific analysis, and financial modeling. For example, Github Copilot, an AI tool that generates code suggestions for programmers, has achieved impressive performance on technical coding questions and, if asked, can provide natural language explanations of how the code it produces works \citep{nguyen2022a, zhao2023}. In addition, beyond learning to predict good outcomes from human-generated data, ML models can implicitly identify characteristics or patterns of behavior that distinguish high and low performers. Generative AI could replace lower-skill workers with AI-based tools or they may use them to help lower-skill, less-experienced workers get up to speed more quickly. Together, the rise of generative AI has the potential to significantly alter the established relationships among technology, labor productivity, and economic inequality \citep{whitehouse2022}.

Despite their potential, generative AI tools face significant challenges in real-world applications.  At a technical level, popular LLM-based tools, such as ChatGPT, have been shown to produce false or misleading information unpredictably, raising concerns about their reliability in high-stakes situations.  While LLM models often perform well on specific tasks in the lab \citep{GPT4TechnicalReport2023, peng2023impact, noyExperimentalEvidenceProductivity2023}, the types of problems that workers encounter in real-world settings are likely to be broader and less predictable.  Furthermore, generative AI tools often require prompts from human operators, yet finding ways to effectively combine human and AI expertise is difficult: for instance, earlier research indicates that decision-support systems integrating AI with human judgment often perform worse than those that rely on humans or AI alone \citep{vaccaro2024combinationshumansaiuseful}. These challenges raise concerns about the ability of AI systems to provide accurate assistance in every circumstance and---perhaps more importantly---workers' capacity to distinguish cases where AI suggestions are effective from those where they are not.  

Finally, the efficacy of new technologies is likely to depend on how they interact with existing workplace structures.  Promising technologies may have more limited effects in practice due to the need for complementary organizational investments, skill development, or business process redesign \cite{brynjolfsson2021productivity}. Because generative AI technologies are only beginning to be used in the workplace, little is currently known about their impacts.  

\section{Our Setting: LLMs for Customer Support}

\subsection{Customer Support and Generative AI}

We study the impact of generative AI in the customer service industry, an industry at the forefront of AI adoption \citep{chui_global_2021}.  Client interactions play a crucial role in building strong customer relationships and company reputation. However, as in many occupations, customer service workers vary widely in productivity \citep{berg2018, syverson2011}.

Newer workers require significant training and time to become more productive.  Turnover is high: industry estimates suggest that 60 percent of agents in contact centers leave each year, costing firms \$10,000 to \$20,000 per agent in the United States \citep{gretzBoostingContactcenterPerformance2018, buesing_customer_2020}. Consequently, the average supervisor spends a large share of their time coaching new agents \citep{berg2018}.  

Customer service is also a setting where there is high variability in the abilities of individual agents. For example, top-performing customer support agents are often more effective at diagnosing the underlying technical issue given a customer's problem description.  They ask more questions before offering a solution, spending more time initially to avoid wasting time later solving the wrong problem. Faced with variable productivity, high turnover, and high training costs, firms are increasingly turning to AI tools that might pick up some of these best practices of top performers \citep{chui_global_2021}.

% Trained on many examples of conversations, generative AI models could pick up some of these best practices. 

At a technical level, customer support is well-suited for current generative AI tools.  Customer-agent conversations can be thought of as a series of pattern-matching problems in which one is looking for a superior sequence of actions.  When confronted with an issue such as ``I can't log in,'' an AI/agent must identify the likely problems and their solutions (``Can you check that caps lock is not on?'').  At the same time, the agent must be attuned to the customer's emotional response, using reassuring rather than patronizing language (``That wasn't stupid of you at all!  I always forget to check that too!''). Because customer service conversations are widely recorded and digitized, pre-trained LLMs can be fine-tuned with examples of both successfully and unsuccessfully resolved conversations. 

\subsection{Data Firm Background} 

We work with a company that provides AI-based customer service support software (hereafter, the ``AI firm'') to study the deployment of its tool at a client firms (hereafter, the ``data firm''). Our data firm is a Fortune 500 company that specializes in business process software for small and medium-sized businesses in the United States.  It employs a variety of chat-based technical support agents, directly and through third-party outsourcing firms.  The majority of agents in our sample work from offices located in the Philippines, with smaller groups working in the United States and  other countries.  Across locations, agents are engaged in a fairly uniform job: answering technical support questions from US-based small business owners.

Customer chats are assigned on the basis of agent availability, with no additional pre-screening.  The questions are often complex and support sessions average 40 minutes, with the majority of the chat spent trying to diagnose the underlying technical problem. The agent's job requires a combination of detailed product knowledge, problem-solving skills, and the ability to handle frustrated customers. While our data firm employed other groups of agents to provide chat-based support for different customer segments---such as self-employed individuals or larger businesses---there was no additional sorting for queries related to US-based small businesses.

Our firm measures productivity using three metrics that are standard in the customer service industry: ``average handle time'' (AHT), the average time an agent takes to finish a chat; ``resolution rate'' (RR), the share of conversations successfully resolved; and ``net promoter score'' (NPS), based on a random post-chat survey that measures customer satisfaction by subtracting the percentage of clients who would not recommend an agent from the percentage who would. A productive agent fields customer chats quickly, while maintaining a high-resolution rate and net promoter score. 

Across locations, agents are organized into teams with a manager who provides feedback and training to agents. Once a week, managers hold one-on-one feedback sessions with each agent. For example, a manager might share the solution to a new software bug, explain the implication of a tax change, or suggest how to better manage customer frustration with technical issues. Agents work individually and the quality of their output does not directly affect others. 

Agents employed by our data firm are generally paid a baseline hourly wage and receive bonuses for hitting specific performance targets, such as for chats per hour or resolution rate. While we lack data on individual pay, the managers we interviewed estimated that performance bonuses accounted for 20 percent to 40 percent of a typical agent's total take-home pay.


\subsection{AI System Design}

The AI system we study is designed to identify conversational patterns that predict efficient call resolution. The system builds on GPT-3 and is fine-tuned on a large dataset of customer-agent conversations labeled with various outcomes, such as call resolution success and handling time.  Our AI firm also up-weights the value of training chats if the chat was conducted by a top performer when training the AI.  Many aspects of successful agent behavior are difficult to quantify, including when to ask clarifying questions, being attentive to customer concerns, deescalating tense situations, adapting communication styles, and explaining complex topics in simple terms.  Explicitly training the AI model on text from top performers helps the AI system pick up on these subtleties in behavior and tone. The AI firm further trains its model using a process similar in spirit to \citet{ouyangTrainingLanguageModels2022} to prioritize agent responses that express empathy, provide appropriate technical documentation, and limit unprofessional language.  This additional training mitigates the potential for hallucinations, and inappropriate responses while helping the LLM distinguish successful behaviors of the top performers, including those they tacitly apply. 

Once deployed, the AI system generates two main types of output: 1) real-time suggestions for how agents should respond to customers and 2) links to the data firm's internal documentation for relevant technical issues. Appendix Figure \ref{afig:AI-sample} illustrates an example of AI assistance. In the chat window (Panel A), Alex, the customer, describes his problem to the agent.  Here, the AI assistant generates two suggested responses (Panel B). The tool has learned that phrases like ``I can definitely assist you with this!'' and ``Happy to help you get this fixed asap'' are associated with positive outcomes. Panel C of Appendix Figure \ref{afig:AI-sample} shows a technical recommendation from the AI system, namely a link to the data firm's internal technical documentation.

% .\footnote{For example, the correct response when a customer says ``I can't track my employee's hours during business trips'' depends on what version of the data firm's software the customer uses. Suppose that the customer has previously mentioned that they are using the premium version. In that case, they should have access to remote mobile device timekeeping, meaning that the support agents need to diagnose and resolve a technical issue that prevents the software from working. If, however, the customer stated that they are using the standard version, then the correct solution is for the customer to upgrade to the premium version in order to access this feature.  For more on context tracking, see, for instance, \citet{dunn2021}.} 

Importantly, the AI system we study is designed to augment, rather than replace, human agents. The output is shown only to the agent, who has full discretion over which, if any, AI suggestions to accept. Giving the agent final authority reduces the likelihood that customers receive off-topic or incorrect responses.  The system does not provide suggestions when it has insufficient training data for a topic, leaving agents to respond on their own. 


\section{Deployment, Data, and Empirical Strategy}

\subsection{AI Rollout}\label{sec:rollout}  
AI assistance was introduced after a randomized control trial (RCT) involving a small number of agents. Appendix Figure \ref{afig:deploymenttimeline} illustrates the timing of the rollout, which primarily took place during the fall of 2020 and winter of 2021. Implementation varied among sites because of limited training resources and the firm's overall budget for AI assistance.

Agents gained access to the AI tool after completing a three-hour online  session conducted by the AI firm. To maintain quality and consistency, training sessions were kept small and exclusively led by one of two employees from the AI firm, both of whom had prior contact center experience and deep knowledge of the AI system. Since they had other full-time responsibilities, the trainers had to limit the number of sessions they could conduct each week. The timing of sessions was adjusted to accommodate the time zones of the data firm's global workforce.

Additionally, because generative AI was costly and relatively untested at that time, the data firm allocated a limited budget for AI deployment. Once the total number of on-boarded agents reached the predefined contractual limit, on-boarding ceased. However, when AI-enabled agents left, their slots were filled by new agents. The capacity and training session scheduling constraints created the variation in AI adoption that we analyze in our study.

Managers at each office oversaw the selection and allocation of agents to training sessions.  In interviews, employees of our AI firm reported that managers sought to minimize customer service disruptions by assigning workers on the same team to different training sessions.  After their initial onboarding session, workers received no additional training on using the AI software. At this time, the AI firm's small product management team did not have the capacity to provide ongoing support to the thousands of agents using the tool.

As a result of these considerations, our AI rollout effectively occurred at the individual level. Within the same team and same office, individuals would be on-boarded to AI assistance at different times. In October 2020, within a team, the average share of active workers with access to AI assistance was only 5 percent, growing to 70 percent in January 2021.  While our analysis primarily focuses on the individual adoption dates, we also provide results in Appendix Table \ref{tab:ivdd_companylocteam} that instrument individual adoption dates with team-level adoption patterns.


\subsection{Summary Statistics}\label{sec:sumstats}

Table \ref{sumstat_agents} provides details on the sample characteristics, divided into four groups: all agents (`all''); agents who never have access to the AI tool during our sample period (``never treated''); pre-AI observations for those who eventually get access (``treated, pre'');and post-AI observations (``treated, post'').  In total, we observe the conversation text and outcomes associated with 3 million chats by 5,172 agents.  Within this, we observe 1.2 million chats by 1,636 agents in the post-AI period.  Most of the agents in our sample, 89 percent, are located outside the United States, mainly in the Philippines.  For each agent, we observe their assigned manager, tenure, geographic location, and employer (the data firm or a subcontractor). 

We rely on several key performance indicators, or outcome variables, all aggregated at the agent-month level, the most granular level with complete data.  Our primary productivity measure is resolutions per hour (RPH), which reflects the number of chats a worker successfully handles per hour. RPH is influenced by several factors, which we also measure individually: the average time handling an individual chat (AHT); the number of chats an agent handles per hour, which accounts for multitasking (CPH); and resolution rate (RR), the share of chats that are successfully resolved. We measure customer satisfaction using the net promoter score from post-call surveys.  While our main outcome measures are at the agent-month level, some data, like chat duration, are available at a more granular chat level.  We construct additional measures of sentiment, topics, and language fluency from chat text.

Our dataset includes average handle time and chats per hour for all agents in our sample. However, subcontractors fail to consistently collect call quality metrics for all agents. As a result, we only observe our omnibus productivity measure---resolutions per hour---for this smaller subset of agents with call quality outcomes. Workers may work only for a portion of the year or part-time, so we calculate year-month observations based solely on the periods that an agent is assigned to chats.
Appendix Section \ref{asec:key_vars} includes a more extensive discussion of our sample construction and key variables. 

Figure \ref{prod-distributions} plots the raw distributions of our outcomes for each of the never, pre-, and post-treatment subgroups.  Several of our main results are readily visible in these raw data.  In Panels A through D, we see that post-treatment agents do better along a range of outcomes, relative to both never-treated agents and pre-treatment agents.  In Panel E, we see no significant differences in surveyed customer satisfaction among treated and non-treated groups. 

Focusing on our main productivity measure, Panel A of Figure \ref{prod-distributions} and Table \ref{sumstat_agents} show never-treated agents resolve an average of 1.7 chats per hour, whereas post-treatment agents resolve 2.5 chats per hour.  Some of this difference may be due to initial selection: treated agents already had higher resolutions per hour prior to AI model deployment (2.0 chats) relative to never treated agents (1.7).  This same pattern appears for chats per hour (Panel C) and resolution rates (Panel D): while ever-treated agents appear to be stronger performers at the outset than agents who are never treated, post-treatment agents perform substantially better.  Looking instead at average handle times (Panel B), we see a starker pattern: pre-treatment and never-treated agents have similar distributions of average handle times, centered at 40 minutes, but post-treatment agents have a lower average handle time of 35 minutes.

These figures, of course, reflect raw differences that do not account for potential confounding factors such as differences in agent experience or differences in selection into treatment.  In the next section, we will more precisely attribute these raw differences to the impact of AI model deployment.

\subsection{Empirical Strategy}

We isolate the causal impact of access to AI recommendations using a standard difference-in-differences regression:
\begin{equation} \label{eq:dd_main}
y_{it}= \delta_{t} + \alpha_{i} + \beta AI_{it} + \gamma X_{it}+ \epsilon_{it}
\end{equation}

Our outcome variables, $y_{it}$, is performance measures for agent $i$ in year-month $t$, with resolutions per hour as our main measure of productivity.  We measure these outcomes in levels, and report percentage changes off the baseline pre-treatment means.  Our main variable of interest is $AI_{it}$, an indicator that equals one if AI assistance is activated for agent $i$ at time $t$.  All regressions include year-month fixed effects, $\delta_{t}$, to control for common, time-varying factors such as tax season or the end of the business quarter.  In our preferred specifications, we also include agent fixed effects, $\alpha_{i}$ , to control for time-invariant differences in productivity across agents and time-varying tenure controls $X_{it}$ (specifically, fixed effects for agent tenure in months).  In our main specifications, we weight each agent-month equally and cluster standard errors at the agent level to reflect that AI access is rolled out individually, but Appendix Tables \ref{tab:robustness_clustering} and \ref{tab:robustdd_nchat_weight} show our results are robust to alternative weightings and clustering. 

A rapidly growing literature has shown that two-way fixed effects regressions deliver consistent estimates only with strong assumptions about the homogeneity of treatment effects and may be biased when treatment effects vary over time or by adoption cohort \citep{cengiz_effect_2019,dechaisemartin2020,sun_estimating_2020,goodman-bacon_2021,callaway_2021, borusyak2022revisiting}. For example, workers may take time to adjust to using the AI system, in which case its impact in the first month may be smaller.  Alternatively, the on-boarding of later cohorts of agents may be smoother, so that their treatment effects may be larger.

We study the dynamics of treatment effects using the interaction-weighted (IW) estimator proposed in \citet{sun_estimating_2020}.  \citet{sun_estimating_2020} show this estimator is consistent assuming parallel trends, no anticipatory behavior, and cohort-specific treatment effects that follow the same dynamic profile.
%\footnote{This last assumption means that treatment effects are allowed to vary over event-time and that average treatment effects can vary across adoption-cohorts (because even if they follow the same event-time profile, we observe different cohorts for different periods of event-time).} 
Both our main differences-in-differences and event study estimates are similar using robust estimators introduced in \citet{dechaisemartin2020}, 
\citet{sun_estimating_2020},
\citet{callaway_2021},
and
\citet{borusyak2022revisiting}, as well as using traditional two-way fixed effects OLS. In fact, Appendix Figure \ref{fig_aiym} shows similar treatment effects across adoption cohorts (e.g. those that received AI access earlier or later, and were thus subject to potentially different versions of the model).  We also show our results are similar when clustering at different levels of granularity (Appendix Table \ref{tab:robustness_clustering}) and instrumenting agent adoption with the date the first worker within the agent's team received AI access (Appendix Table \ref{tab:ivdd_companylocteam}). 

\section{Main Results}

\subsection{Overall Impacts}
Table \ref{tab:dd_main} examines the impact of the deployment of the AI model on our primary measure of productivity, resolutions per hour, using a standard two-way fixed effects model.  In Column 1, we show, controlling for time and location fixed effects, access to AI recommendations increases resolutions per hour by 0.47 chats, up 23.9 percent from their pre-treatment mean of 1.97.  In Column 2, we include fixed effects for individual agents to account for potential differences between treated and untreated agents.  In Column 3, we include additional fixed effects for agent tenure in months to account for time-varying experience levels.  As we add controls, our effects fall slightly, so that, with agent and tenure fixed effects, we find the deployment of AI increases RPH by 0.30 chats or 15.2 percent. 

Figure \ref{fig:es_main_AS} shows the accompanying IW event study estimates of \citet{sun_estimating_2020} for the impact of AI assistance on RPH.  We find a substantial and immediate increase in productivity in the first month of deployment.  This effect grows slightly in the second month and remains stable and persistent up to the end of our sample. 

In Table \ref{tab:dd_oth}, we report additional results using our preferred specification with fixed effects for year-month, agent, location, and agent tenure.  Column 1 documents a 3.7 minute decrease in the average duration of customer chats, an 8.5 percent decline from the baseline mean of 43 minutes (shorter handle times are considered better).  Next, Column 2 indicates a 0.37 unit increase in the number of chats that an agent can handle per hour.  Relative to a baseline mean of 2.4, this represents an increase of roughly 15 percent.  Unlike average handle time, chats per hour account for the possibility that agents may handle multiple chats simultaneously.  The fact that we find a stronger effect on this outcome suggests that AI enables agents to both speed up chats and multitask more effectively.

Column 3 of Table \ref{tab:dd_oth} indicates a small 1.3 percentage point increase in chat resolution rates.  This effect is economically modest and insignificant given a high baseline resolution rate of 82 percent. We interpret this as evidence that improvements in chat handling do not come at the expense of problem solving on average.  Finally, Column 4 finds no economically significant change in customer satisfaction, as measured by net promoter scores: the coefficient is -0.12 percentage points and the mean is 80 percent. 

Appendix Figure \ref{afig:es_oth} presents the accompanying event studies for additional outcomes.  We see immediate impacts on average handle time (Panel A) and chats per hour (Panel B), and relatively flat patterns for resolution rate (Panel C) and customer satisfaction (Panel D).  We interpret these findings as saying that, on average, AI assistance increases productivity without negatively impacting resolution rates and surveyed customer satisfaction. 

\subsubsection{Randomized Control Trial Analysis} \label{sec:rct}

In August 2020, our data firm conducted a pilot analysis involving approximately 50 workers, with about half randomized into treatment. Unfortunately, we do not have information on the control group workers that were part of the experiment, so we compare our treated group to all remaining untreated agents. Analysis of the randomized control trial, in Appendix Table \ref{rctdd.tex}, shows similar effects on productivity as in our main sample. Appendix Figure \ref{fig_rct} reports the accompanying event studies for our various outcomes, which are also similar to our main results.  


\subsubsection{Robustness} \label{sec:instrumentingadoption}
%\subsubsection{Instrumenting Individual AI adoption} 

During the roll-out process, managers decided which agents to onboard onto the AI system and scheduled when their training would occur. If managers allocated AI access to stronger workers first, our OLS results could overstate the impacts of AI.  To address this, in Appendix Table \ref{tab:ivdd_companylocteam}, we instrument an individual agent's AI adoption date with the first adoption date of the worker's company,  office location, and team. The effects on average handle time and chats per hour are essentially identical to those under our main specification. However, resolutions increase by 0.55 chats per hour, compared with 0.30 in our main finding. We attribute the larger impact on resolutions per hour to the fact that this IV approach estimates a significant and larger impact on resolution rates.

% \footnote{Information on team is missing for some workers.  Workers without this information are grouped into a single ``missing'' team.} 
% Although AI-access was in theory introduced across the entire firm, variation in training schedules and time zones made it so that some teams began the process sooner than others.  Using this approach, w

% Unlike traditional OLS, these estimators avoid comparing between newly treated and already treated units.
We also show similar results using alternative estimators and at different levels of clustering and weighting. Appendix Table \ref{atab:dd_main_robust} finds similar results using alternative difference-in-difference estimators introduced in \citet{dechaisemartin2020}, \citet{sun_estimating_2020}, \citet{callaway_2021}, and \citet{borusyak2022revisiting}. Similarly, Appendix Figure \ref{afig:es_main_alt} reports that our results are similar under alternative event study estimators: \citet{dechaisemartin2020}, \citet{callaway_2021}, \citet{borusyak2022revisiting}, and traditional two-way fixed effects.  Next, in Appendix Table \ref{tab:robustness_clustering}, we show our standard errors are similar whether clustering at the individual level, team level, or geographic location level.  Finally, we explore robustness to alternative weighting.  In Appendix Table \ref{tab:robustdd_nchat_weight}, we weight agent-month observations by the number of customer chats a worker conducts. Re-weighting generates similar results to our main, equally-weighted, specifications. 
%While our standard errors are smaller when clustering at the agent level, our results are still highly significant when clustering at agents' geographic location, the most conservative level.
%\footnote{Geographic location is usually available at the city level (e.g. Reno, Nevada) and one geographic location may contain firms employed by different subcontractors, although a small share workers work remotely.}  

  

\subsection{Heterogeneity by Agent Skill and Tenure}

There is substantial interest in the distributional consequences of AI-based technologies. An extensive literature suggests that earlier waves of information technology complemented high-skill workers, with effects on productivity, labor demand and wage differentials. Together with important changes in relative supply and demand for skilled labor, these changes shaped patterns of wage inequality in the labor market \citep{katzgoldinorigins1998, katzgoldineducationtech2008}.  Unlike earlier waves of IT, however, generative AI does not simply execute routine tasks. Instead, as outlined in Section \ref{sec:aibackground}, AI models identify patterns in data that replicate the behaviors of many types of workers, including those engaged in non-routine, creative, or knowledge-based tasks.  These fundamental technical differences suggest that generative AI may impact different workers in different ways. 

\subsubsection{Pre-treatment Worker Skill}
We explore two components that are important for understanding the distributional consequences of AI adoption: its impacts by agent productivity and tenure. We measure an agent's ``skill'' using an index incorporating three key performance indicators: call handling speed, issue resolution rates, and customer satisfaction.  To construct this index, we compute an agent's ranking within its employer company-month for each component productivity measure and then average these rankings into a single index.  Next, we calculate the average index value over the three months for each agent, to smooth out month-to-month shocks in agent performance.  An agent in the top quintile of this productivity index demonstrates excellence across all three metrics: efficient call handling, high issue resolution rates, and superior customer satisfaction scores. 

Panel A of Figure \ref{fig:dd_workerhet_skill} shows how our productivity effects vary  across workers in each quintile of our skill index, measured in the month prior to AI access.  To isolate the impact of worker skill, our regression specification, available in Appendix Section \ref{asec:specifications}, includes a set of fixed effects for months of worker tenure. 
% pre treatment rph means by skill bin: 1 1.48, 2 - 1.91, 3 -1.99, 4 - 2.15, 5 - 2.47
% event study coeffs by skill bin: 1 .53, 2 .32, 3 .30, 4. 26. 5 .02
% percentage changes = 36, 17, 15, 12, .01
We find the productivity impact of AI assistance is most pronounced for workers in the lowest skill quintile (leftmost side), who see an increase of 0.5 in resolutions per hour, or 36 percent. In contrast, AI assistance does not lead to any significant change in productivity for the most skilled workers (rightmost side).

In Panels B through E of Figure \ref{fig:dd_workerhet_skill} we show that less-skilled agents consistently see the largest gains across our other outcomes as well.  For the highest-skilled workers, we find mixed results: a zero effect on average handle time (Panel B); a small but positive effect for chats per hour (Panel C); and, interestingly, small but statistically significant \textit{decreases} in resolution rates and customer satisfaction (Panels D and E).  

These results are consistent with the idea that generative AI tools may function by exposing lower-skill workers to the best practices of higher-skill workers.  Lower-skill workers benefit because AI assistance provides them with new solutions, whereas the best performers may see little benefit from being exposed to their own best practices.  Indeed, the negative effects along measures of chat quality---resolution rate and customer satisfaction---suggests that AI recommendations may distract top performers, or lead them to choose the faster or less cognitively taxing option (following suggestions) rather than taking the time to come up with their own responses.  Addressing this outcome is potentially important because the conversations of top agents are used for ongoing AI training.

Our results could be driven by mean reversion: Agents who performed well just prior to AI-adoption may see a natural decline in their productivity afterward, while lower-performing agents may rebound.  To address this concern, we plot raw resolutions per hour in event-time, graphed by skill tercile at AI treatment in Appendix Figure \ref{afig:meanrev}. If mean reversion were driving our observed effects, we would expect to see a convergence of productivity levels after treatment, with top-tercile agents showing decreased performance and the least-skilled agents demonstrating improved output. However, our analysis reveals a consistent linear increase in productivity across all skill levels after AI implementation, with no strong evidence of mean reversion, suggesting that productivity gains are attributable to AI assistance.

\subsubsection{Pre-treatment Worker Experience}
% pre treatment rph means by tenure bin: 1 1.44, 2 - 1.74, 3 -1.87, 4 - 1.94, 5 - 2.18
% event study coeffs by tenure bin: 1 .70, 2 .39, 3 .36, 4. 33. 5 .02
% percentage changes = 49, 22, 19, 17, .01
Next, we repeat our previous analysis for agent tenure to understand how the treatment effects of AI access vary by worker experience. To do so, we divide agents into five groups based on their months of tenure at the time the AI model is introduced.  Some agents have less than a month of tenure when they receive AI access, while others have more than a year of experience.  To isolate the impact of worker tenure, this analysis controls for the worker skill quintile at AI adoption, with the regression specification located in Appendix Section \ref{asec:specifications}. 

In Panel A of Figure \ref{fig:dd_workerhet_tenure}, we see a clear monotonic pattern in which the least experienced agents see the greatest gains in resolutions per hour.  Agents with less than 1 month of tenure improve by 0.7 resolutions per hour, with larger effects for less-experienced workers.  In contrast, we see no effect for agents with more than a year of tenure.  

In Panels B through E of Figure \ref{fig:dd_workerhet_tenure}, we report results for other outcomes.  In Panels B and C, we see that AI assistance generates large gains in call handling efficiency, measured by average handle times and chats per hour, respectively, among the newest workers.  In Panels D and E, we find positive impacts of AI assistance on chat quality, as measured by resolution rates and customer satisfaction, respectively.  For the most experienced workers, we see modest positive effects for average handle time (Panel B), positive but statistically insignificant effects on chats per hour (Panel C), and small but statistically significant negative effects for measures of call quality and customer satisfaction (Panels D and E). 

Overall, these patterns are very similar to our findings for agent skill, even though our regressions are designed to isolate the distinct roles of skill and experience -- our skill regressions control for experience and vice versa.  This suggests that even within the same task, access to AI systems disproportionately improves the performance of both novice and less skilled workers.

\subsubsection{Moving Down the Experience Curve}

To further explore how AI assistance impacts newer workers, we examine how worker productivity evolves on the job.
% \footnote{We avoid the term ``learning curve'' because we cannot distinguish if workers are learning or merely following recommendations.}  
In Figure \ref{fig:learning_bycohort}, we plot productivity variables by agent tenure for three distinct groups: agents who never receive access to the AI model (``never treated''), those who have access from the time they join the firm (``always treated''), and those who receive access in their fifth month with the firm (``treated 5 mo.'').  

We see that all agents begin around at 1.8 resolutions per hour.  Never-treated workers (long dashed blue line) slowly improve their productivity with experience, reaching approximately 2.5 resolutions per hour 8 to 10 months later.  In contrast, workers who always have access to AI assistance (short dashed red line) increase their productivity to 2.5 resolutions per hour after only two months and continue to improve until they are resolving more than 3 chats per hour after five months of tenure.%\footnote{Our sample ends here because we have very few observations more than five months after treatment.} 
Comparing just these two groups suggests that access to AI recommendations helps workers move more quickly down the experience curve and reduces ramp-up time. 

The final group in Panel A tracks workers who begin their tenure without access to AI assistance but who receive access after five months on the job (solid green line).  These workers initially improve at the same rate as never-treated workers, but after gaining AI access in month 5, their productivity begins to more rapidly increase, following the same trajectory as the always-treated agents. These findings demonstrate AI assistance not only accelerates ramp-up for new workers, but also improves the rate at which even experienced workers improve in their roles.

In Appendix Figure \ref{afig:learning_bycohort}, we plot these curves for other outcomes. We see clear evidence that the experience curve for always-treated agents is steeper for handle time, chats per hour, and resolution rates (Panels A through C). Panel D follows a similar but noisier pattern for customer satisfaction. Across many of the outcomes we examine, agents with two months of tenure and access to AI assistance perform as well as or better than agents with more than six months of tenure who do not have access. AI assistance alters the relationship between on-the-job productivity and time, with potential implications for how firms might value prior experience, or approach training and worker development.

\section{Adherence, Learning, Topic Handling, and Conversational Change}

We conduct the following analyses to understand better the mechanisms driving our main results. We examine patterns in how workers of varying skills engage AI recommendations. We look at how exposure to AI helps agents' master their jobs, sharpening their diagnostic skills and language fluency, and how AI assistance influences the communication patterns of higher- and lower-skill workers. 

%Taken together, our results suggest that examining and following AI recommendations helps workers---particularly lower-skilled workers---learn to adopt best practices gathered from higher-skill and more experienced agents.  
% In this section, we conduct a variety of analyses to better understand the mechanisms driving our main results.  

% First, we examine how workers engage with AI recommendations.  We show that workers are selective about the recommendations they adopt, following the recommendations 35\% on average.  We find that the returns to AI assistance are highest for workers who choose to follow the recommendations.  Consistent with a story in which workers find AI recommendations helpful, we show that adherence rates increase over time for all workers, especially among more experienced workers: by the end of our sample, we see similar adherence rates across worker tenure and skill. 

% Second, we explore whether AI-assistance helps workers learn.  Using information on software outages in which AI assistance is temporarily unavailable, we provide evidence that exposure to AI leads to durable changes in worker skills. We find that workers exposed to AI recommendations continue to perform better during outages, and this effect is greater after more exposure and for agents who follow AI recommendations more closely when the software is working.

% Next, we provide evidence on two ways in which AI assistance may improve workers' conversations: providing them with the information to address specific customer problems and refining their overall conversational fluency.  Our findings indicate that AI access is particularly beneficial when workers encounter less routine issues, where they might otherwise struggle to find solutions. In addition, we also show that AI assistance improves workers' English language skills, with the most substantial gains seen among those with the lowest initial proficiency. 

% Lastly, we analyze the text of conversations to assess how AI assistance influences the communication patterns of higher- and lower-skill workers.  In particular, we track how an agent's words change before and after AI adoption and find evidence of larger textual shifts for lower-skill workers.  We also show that the similarity of conversations conducted by high- and low-skill agents increases following AI-adoption.  These results suggest that AI assistance may help lower-skill agents converge towards their higher-skill peers.

% Taken together, our results suggest that examining and following AI recommendations helps workers---particularly lower-skilled workers---learn to adopt best practices gathered from higher-skill and more experienced agents.  


\subsection{Adherence to AI Recommendations} \label{sec:receptivity}

The AI tool we study makes suggestions, but agents are ultimately responsible for what they say to the customer. Thus far, our analysis evaluates the effect of AI assistance, irrespective of the frequency with which users adhere to its suggestions.
%In our main results, we estimate how access to the AI tool impacts outcomes regardless of how frequently agents follow its recommendations.  
Here, we examine how closely agents adhere to AI recommendations and document the association between adherence and returns to adoption.  
%We measure ``adherence'' starting at the chat level, using the share of AI recommendations that each agent follows.  Our AI firm codes agents as having adhered to a recommendation if they either click to copy the suggested AI text or if they self-input something very similar.  We take  this chat-level measure and aggregate it to the agent-month level.  
We define ``adherence'' as the proportion of AI suggestions an agent typically adopts, when an AI suggestion is generated, omitting messages where the AI does not make suggestions. The AI company considers an agent to have adhered when they either directly copy the AI's proposed text or manually enter highly similar content. To gauge initial adherence, we classify each treated agent into a quintile based on their level of adherence during their first month using the AI tool. 

%so agents should ignore these recommendations.   

Panel A of Figure \ref{fig:dd_byreceptivity} shows the distribution of average agent-month-level adherence for our post-AI sample, weighted by the log number of AI recommendations provided to that agent in that month.  The average adherence rate is 38 percent, with an interquartile range of 23 percent to 50 percent: Agents frequently ignore recommendations.  In fact, the share of recommendations followed is similar to the share of other publicly reported numbers for generative AI tools; a study of GitHub Copilot reports that individual developers use 27 percent to 46 percent of code recommendations \citep{zhao2023}. Such behavior may be appropriate, given that AI models may make incorrect or irrelevant suggestions.  In Appendix Figure \ref{fig:adherence_decomposed}, we further show that the variation in adherence is similar within locations and teams, indicating that it is not driven by some organizational segments being systematically more supportive than others
%so agents should ignore these recommendations.   

Panel B of Figure \ref{fig:dd_byreceptivity} shows that \textit{returns} to AI model deployment are higher when agents follow recommendations.  We measure this by dividing agents into quintiles based on the share of AI recommendations they follow in the first month of AI access. Following Equation \ref{aeq:dd_adherence}, we separately estimate the impact of AI assistance for each group, including year-month, agent, and agent tenure fixed effects.

We find a steady and monotonic increase in returns by agent adherence: Among agents in the lowest quintile, we still see a 10 percent gain in productivity, but for agents in the highest quintile, the estimated impact is over twice as high, close to 25 percent.  Appendix Figure \ref{afig:dd_byreceptivity} shows the results for our other four outcome measures. The positive correlation between adherence and returns holds most strongly for average handle time (Panel A) and chats per hour (Panel B), and more noisily for resolution rate (Panel C) and customer satisfaction (Panel D). 

Our results are consistent with the idea that there is a treatment effect of following AI recommendations on productivity.  We note, however, that this relationship could also be driven by other factors: selection (agents who choose to adhere are more productive for other reasons); or selection on gains (agents who follow recommendations are those with the greatest returns).  We then consider the worker's revealed preference: Do they continue to follow AI recommendations over time?  If our results were driven purely by selection, we would expect workers with low adherence to continue having low adherence, since it was optimal for them to do so.

Appendix Figure \ref{fig:receptivity_overtime} plots the evolution of AI adherence over time, for various categories of agents.  Panel A begins by considering agents who differ in their initial AI compliance, which we categorize based on terciles of AI adherence in the first month of model deployment (``initial adherence'').  Here, we see that compliance either stays stable or grows over time.  The most initially compliant agents continue to comply at the same rates (just above 50 percent).  The least initially compliant agents increase their compliance over time: Those in the bottom tercile initially follow recommendations less than 20 percent of the time, but by month five their compliance rates have increased by over 50 percent.  

Next, Panel B divides workers up by tenure at the time of AI deployment.  More senior workers are initially less likely to follow AI recommendations: 30 percent for those with more than a year of tenure compared with 37 percent for those with less than three months of tenure.  However, over time, all workers increase adherence, with more senior workers doing so faster, and the groups converge five months after deployment.  

In Panel C, we show the same analysis by worker skill at AI deployment.  Here, we see that compliance rates are similar across skill groups, and all groups increase their compliance over time.  In Appendix Figure \ref{fig:receptivity_overtime_agentfe} we show these patterns are robust to examining within-agent changes in adherence (that is, adherence rates residualized by agent fixed effects), suggesting that increases in adherence over time are not driven exclusively by less adherent agents leaving.     

The results in Figures \ref{fig:receptivity_overtime} and \ref{fig:receptivity_overtime_agentfe} are consistent with agents, particularly those who are initially more skeptical, coming to value AI recommendations over time.  We note, however, that high-skill agents increase their adherence as quickly as their lower-skill peers, even though their productivity gains are smaller and---in the case of some quality measures---even negative.  This suggests an alternate possibility: some agents may be over-relying on AI recommendations beyond what is optimal in the long run. Top agents, in particular, may see little additional value in taking the time to provide the highest quality inputs when an adequate AI suggestion is readily available.  High AI adherence in the present may then reduce the quality or diversity of solutions used for AI training in the future.  However, in the short run, our analysis finds no evidence that the model is declining in quality over our sample period. In Appendix Figure \ref{fig_aiym}, we show that workers who received later access to the AI system---and therefore to a more recently updated version---had similar first-month treatment effects as those who received access to an earlier version of the model. 


\subsection{Worker Learning} \label{sec:learning}

A key question raised by our findings so far is whether these improvements in productivity and changes in communication patterns reflect durable changes in the human capital of workers or simply their growing reliance on AI assistance.  In the latter case, the introduction of AI assistance could actually lead to an erosion in human capital, and we would expect treated workers to be less able to address customer questions if they are no longer able to access AI assistance. For example, research in cognitive science has shown that individuals learn less about spatial navigation when they follow GPS directions, relative to using a traditional map \citep{Brügger2019}.

We examine how workers perform during periods in which they are not able to access AI recommendations due to technical issues at the AI firm.  Outages occur occasionally in our data and can last anywhere from a few minutes to a few hours.  During an outage, the system fails to provide recommendations to some, but not necessarily all, workers.  For example, outages may affect agents who log into their computers after the system crashes, but not agents working at the same time who had signed in earlier. Outages may also affect workers using one physical server but not another. Our AI firm tracks the most significant outages in order to perform technical reviews of what went wrong.  We compile these system reports to identify periods in which a significant fraction of chats are affected by outages.

Appendix Figure \ref{afig:outage_example} shows an example of such an outage, which occurred on September 10, 2020.  The $y$-axis plots the share of post-treatment chats (e.g. those occurring after the AI system has been deployed for a given agent) for which the AI software does not provide any suggestions, aggregated to the hour level.  The $x$-axis tracks hours in days leading up to and following the outage event (hours with fewer than 15 post-treatment chats are plotted as zeros for figure clarity).  During non-outage periods, the share of chats without AI recommendations is typically 30 to 40 percent, reflecting that the AI system does not normally generate recommendations in response to all messages.  On the morning of September 10th, however, we see a notable spike in the number of chats without recommendations, increasing to almost 100 percent. Records from our AI firm indicate that this outage was caused by a software engineer running a load test that crashed the system.   

Figure \ref{fig:ES_outage} examines the impact of access to the AI system for chats that occur during and outside these outage periods.  These regressions are estimated at the individual chat level in order to precisely compare conversations that occurred during outage periods with those that did not. Because we do not have information on chat resolution at this level of granularity, our main outcome measure is chat duration. Panel A considers the impact of AI assistance using only post-adoption periods in which the AI system is not impacted by a software outage.  Consistent with our main results, we see an immediate decline in the duration of individual chats by approximately 10 percent to 15 percent.  

In Panel B, we use the same pre-treatment observations, but now restrict to post-adoption periods that are impacted by large outages. We find that even during outage periods when the AI system is not working, AI-exposed agents continue to handle calls faster (equivalent to 15 percent to 25 percent declines in chat duration).  Because AI outages are rare, our estimates are noisy, and could reflect differences in the types of chats that are seen during outage periods than during non-outage periods. However, focusing on the size of estimated effects over time, an interesting pattern emerges.  Rather than declining immediately post-adoption and staying largely stable as we see in Panel A for non-outage periods, Panel B shows that the benefit of exposure to AI assistance increases with time during outage periods.  That is, if an outage occurs one month after AI adoption, workers do not handle the chat much more quickly than their pre-adoption baseline.  Yet, if an outage occurs after three months of exposure to AI recommendations, workers handle the chat faster---even though they are not receiving direct AI assistance in either case.  

Panel B highlights the potential scope for improving existing employee training practices. Prior to AI assistance, training was limited to brief weekly coaching sessions where managers reviewed select conversations and provided feedback. However, by necessity, managers can only provide feedback on a small fraction of the conversations an agent conducts.  Moreover, because managers are often short on time and may lack training, they often simply point out weak metrics (``you need to reduce your handling time'') rather than identifying strategies for how an agent could better approach a problem (``you need to ask more questions at the beginning to diagnose the issue better.'') Such coaching can be ineffective and can be counterproductive to employee engagement \citep{berg2018}. In contrast, AI assistance offers workers specific, real-time, actionable suggestions, potentially addressing a limitation of traditional coaching methods.

To better understand how learning might occur, in Panels C and D of Figure \ref{fig:ES_outage}, we divide our main study of outage events by the initial adherence of the worker to AI, as described in Section \ref{sec:receptivity}. When a worker chooses not to follow a particular AI recommendation, they miss the opportunity to observe how the customer might respond.  AI suggestions may prompt workers to communicate in ways that differ from their natural style, such as by expressing more enthusiasm or empathy, or by frequently pausing to recap the conversation.  Workers who do not try out these recommendations may never realize that customers could react positively to them.

Panel C reveals that workers with high initial adherence to AI recommendations experience significant and rapid declines in chat processing times, even during outages, relative to their pre-adoption baseline. In contrast, Panel D shows no such improvement for workers who frequently deviate from AI suggestions; they see no reduction in chat times during outage periods, even after prolonged AI access.  

These findings suggest that workers learn more by actively engaging with AI suggestions and observing first-hand how customers respond. These findings are consistent with other evidence from education that higher adherence and engagement with LLM-generated responses positively impacted learning \citep{kuman2023}. In addition to directly improving productivity, exposure to AI assistance could supplement existing on-the-job training programs. 

\subsection{Handling Routine and Non-Routine Topics}

In addition to varying by the characteristics of the worker, the impact of AI assistance could depend on the types of problems it is asked to resolve.  Agents encounter customer questions that range from common requests for help onboarding an employee or changing a password to less common issues such as setting up wage garnishments in child support cases or ensuring compliance with international tax treaties. We examine how the impact of AI-assistance varies between more and less routine customer problems. We use Gemini, a large language model developed by Google DeepMind, to classify the interactions into topic categories. The details of this process, along with our human validation of the LLM classification process, are described in Appendix Section \ref{asec:key_vars} \citep{geminipaper}.

% Agents in our data encounter a diverse set of customer questions, ranging from common requests for help in onboarding an employee or changing a password, to less common questions such as setting up wage garnishments in child support cases or ensuring compliance with international tax treaties.  When an AI model is trained on these conversations, it will naturally have more examples of the most common problems, raising the question of whether it can effectively provide suggestions for topics it has encountered less frequently.  

Appendix Figure \ref{hist_topic} reports the distribution of conversation topics in our dataset.  Unsurprisingly, we observe a small number of frequent issues, accompanied by a long tail of less common problems. Specifically, the two most prevalent topics---payroll and taxes, and account access and management---comprise half of all conversations and the top 16 topics represent over 90 percent of all chats.  

To evaluate the impact of AI assistance based on the frequency of customer inquiries, we categorize conversations into four distinct groups. The ``Payroll/Account'' category, comprising 50 percent of all chats, includes inquiries related to payroll, taxes, and account access and management. The next 25 percent of chats covers five additional categories, including those dealing with bank transfers or managing subscriptions.  The following 15 percent of chats encompass nine additional topics, while the final 10 percent of chats consists of all the remaining topics. Our regression, in Appendix Section \ref{asec:specifications}, is conducted at the chat level, with a focus on chat duration.

Panel A of Figure \ref{fig:dd_bytopic} shows the average treatment effect of AI assistance based on how common the inquiry is.  The pattern is non-monotonic and suggests that AI-assistance has the greatest impact on workers' ability to handle problems that are moderately rare.  Workers with access to AI assistance handle the most routine problems---payroll and account management---about 4 to 5 minutes faster, which corresponds to an approximately 10 percent decrease from the pre-treatment mean duration for these topics.   We see the largest decline, 5 to 6 minutes, for issues that are in the 75th to 90th percentiles of topic rarity, corresponding to a 14 percent reduction from the pre-treatment means for those topics.  Finally, we see a smaller 4 minute or 11 percent decrease for the most rare problems.  It should be noted that the system does not provide suggestions at all when there was insufficient training data.  

These results highlight the difference between the technical quality of an AI system and its potential productivity impacts in real-world settings. AI models generally perform better when trained on large datasets, which provide diverse examples and richer contextual information. Such datasets enable the model to learn more robust and generalizable patterns while reducing the risk of overfitting \citep{halevy2009}.  Consequently, we might expect an AI system to function best when dealing with routine problems, where training data are abundant. 

However, the value of AI systems is less straightforward when they are used to complement human workers.  Customer service agents, especially those dealing with common issues, are specifically trained to address these routine problems and become most experienced answering them. For example, even novice workers are likely to know how to reset a customer's password.  In such cases, access to even high-quality AI assistance may not have a large complementary impact. Rather, as our findings suggest, the impact of an AI system on workplace productivity depends critically on its capabilities relative to workers' baseline skills. The greatest productivity gains may occur not where the AI system is most capable in absolute terms, but where its capabilities most effectively complement or exceed those of human workers.

In our setting, the heterogeneous impact of AI access appears to reflect both factors. AI access has the smallest reduction in handle time for problems where human agents are already well trained (very routine problems) or where its training data may be sparse (very rare problems).  We see the largest improvements in the handle time for moderately uncommon problems.  The AI system is likely to have enough training data to assess these problems, while individual agents are less likely to have had much first-hand experience.  For example, the AI-recommended links to potentially relevant technical documentation may be particularly valuable for the types of cases where agents otherwise would need to search for an answer. 

To examine the role of agent-specific experience, Panel B of Figure \ref{fig:dd_bytopic} plots the impact of AI assistance on chat duration by quartiles of topic frequency with respect to an individual agent, controlling for the overall frequency of a problem.  AI assistance reduces conversation times 15 percent for the least common problems compared with 10 percent for the most common.  Once we control for a topic's overall frequency, we find a monotonic relationship between agent-specific exposure to a problem and the impact of AI. That is, holding constant the AI model's exposure to a problem, the impact of AI assistance is greatest for problems that a specific agent is least exposed to.  While AI in isolation may be most effective where training data is most plentiful, the marginal value of AI assistance appears to be highest where humans have a greater need for AI input. 


\subsection{Conversational Style}

%Top customer service agents not only resolve technical issues but also communicate in a way that ensures customer satisfaction.  For our US-based customers, this involves clearly writing in fluent English and displaying subtle interpersonal skills. Our analysis provides suggestive evidence that AI adoption helps lower-skill workers communicate more like their higher-skill peers.
%Our previous analysis examined how the effectiveness of AI assistance varies based on the technical nature of the problems being addressed. However, top customer service agents not only resolve technical issues but also communicate in a way that ensures customer satisfaction.  For our US-based customers, this involves clearly writing in fluent English and displaying subtle interpersonal skills.  In this section, we first demonstrate that AI assistance improves overall language fluency.  Then, to capture less quantifiable changes in communication style, we compare the evolution of conversational text produced by different workers.  Our analysis provides suggestive evidence that AI adoption helps lower-skill workers communicate more like their higher-skill peers.


\subsubsection{English Fluency} \label{sec:englishproficiency}

The ability to communicate in clear, idiomatic English is crucial to customer satisfaction and the job performance of contact workers serving US customers. In our data, 80 percent of the agents are based in the Philippines, where many residents are fluent English speakers for various cultural and historical reasons.  However, cultural differences and language nuances occasionally lead to misunderstandings or a sense of disconnect, even when an agent's technical language skills are strong. In this section, we assess how AI assistance influences workers' ability to communicate clearly. 

% \footnote{English is an official language of the Philippines and  Filipinos, especially those who are college educated, often have near-native fluency due to early exposure in schools.  Further, the Philippines is a former American colony and American English tends to be the dominant version of English used in the media.}

We measure the proficiency of text in two ways: its \textit{comprehensiblity} and its \textit{native fluency}. The comprehensibility score assesses whether the agent produces text that is cogent and easy to understand, using a scale of 1 to 5, where 1 indicates ``very difficult to comprehend'' and 5 signifies ``very fluent and easily understandable, with no significant errors.''  In contrast, native fluency focuses on whether the text was likely to have been produced by a native speaker of American English.  Native fluency is based on the Interagency Language Roundtable ``functionally native'' proficiency standard.  The native fluency score is also on a 5 point scale where 1 indicates a writer is ``Definitely not a native American English speaker'' and 5 indicates they definitely are.  For instance, ``I could care less'' is grammatically incorrect, but a common English-language expression. On the other hand, Filipino agents often use the greeting ``to have a blessed day,'' which is grammatically correct, but not a common greeting in the United States.  We use Gemini, an LLM, to score agents' text in each conversation along these two dimensions. For more information on our specific approach, prompts, and validation tests, see Appendix \ref{asec:language_scoring}. 

The general level of both comprehensibility and native fluency is high. Prior to having AI access, the interquartile range of comprehensibility scores was 4.28 to 4.67; for native fluency it was 4.26 to 4.65.  Despite this high baseline level, we find clear evidence that access to AI assistance increases proficiency scores.  Appendix Figure \ref{afig:englishfluency1}  presents the raw pre- and post-AI distribution of comprehensibility and native fluency scores for never-treated workers, pre-treatment workers, and post-treatment workers.  The never treated and pre-treatment workers have identical distributions, but we see markedly higher scores for post-treatment workers.  In Panels A and B of Figure \ref{fig:englishfluency}, we report the accompanying event studies which show AI access leads to large improvements in both comprehensibility and native fluency.  Finally, in Panels C and D of Figure \ref{fig:englishfluency}, we report separate coefficients for US- and Philippines-based workers.  We see a positive impact for all workers, but a larger improvement for workers based in the Philippines. 


\subsubsection{Textual Convergence}

The analysis above focuses on an important, but narrow, aspect of how workers communicate. To gain a broader understanding of AI's influence on communication patterns, we examine how the text produced by workers evolves over time: do they change how they write relative to their pre-AI baseline, and does AI access impact the relative communication patterns of high and low skill workers? Because tacit knowledge is, by definition, not something that can be codified as a set of rules, we examine the overall textual similarity of conversations using textual embeddings, rather than looking for the presence of specific formulaic phrases \citep{sentence2023}.

% within worker standard dev of cosine similarity pre-treatmetn is .10

Panel A of Appendix Figure \ref{afig:text_change} plots the evolution of agents' communication over time, before and after access to AI assistance. We compute the cosine similarity of agents' text in each given event-time week to their own chats from the month before AI deployment (week -4 to week -1). Cosine similarity runs from 0 to 1, with 0 meaning two pieces of text are orthogonal (when represented as semantic vectors), and 1 indicating exact semantic similarity.  

Prior to the deployment of AI, the similarity between a worker's own text from month to month is stable, at 0.67, which reflects consistency in an individual agent's language use, while also capturing differences in the topics and customers that she faces.  However, following AI deployment, the similarity of agents' text drops. The drop is equivalent to about half of a standard deviation of within-agent cosine similarity across the pre-period. This is consistent with the idea that AI assistance changes the content of agents' messages, rather than merely leading workers to type the same content but faster.  Panel B of Appendix Figure \ref{afig:text_change} plots the average change in textual content separately by pre-AI worker skill. Lower-skill agents experience greater textual change after AI adoption, relative to top performers. 

% standard dev of H/L similiary score is .12
%\footnote{For non-AI users, we define skill levels based on monthly quintiles of our skill index. For AI users, we use skill quintiles at the time of AI deployment.} 
We find across-worker changes in communication changes with AI access.  Panel C of Appendix Figure \ref{afig:text_change} plots the cosine similarity between high- and low-skill agents at specific moments in calendar time, separately for workers without (blue dots) and with (red diamonds) access to AI assistance. For non-AI users, we define skill levels based on monthly quintiles of our skill index. For AI users, we use skill quintiles at the time of AI deployment.  Without AI, high- and low-productivity workers show a moderate level of similarity in their language use, with an average cosine similarity between high and low workers of 0.55.  This similarity remains stable over time, suggesting that there are no divergent trends between skill groups that do not have access to AI assistance.  Post-AI adoption, however, text similarity between high- and low-skill workers begins increasing, from 0.55 to 0.61. While this change may seem modest, it represents a substantial narrowing of language gaps, given that the average similarity of a high-skill worker's own pre- and post-AI text is only 0.67. The change is equivalent to half of a standard deviation of the average high and low worker textual similarity.  

Taken together, the patterns in Appendix Figure \ref{afig:text_change} are consistent with AI assistance leading to more pronounced changes in how lower-skill workers communicate and  ultimately to their communicating more like high-skill workers. We caution, however, that changes in agent text can reflect many factors that are not directly related to a worker's style or tacit skills, such as changes in conversation topics driven by customers.  As a result, this analysis is only suggestive. 


\section{Effects on the Experience of Work}

%Qualitative studies suggest that working conditions for contact center agents can be unpleasant.  They are regularly exposed to challenging, emotionally charged and anonymized conversations, and called upon to absorb customer frustrations while limiting one's own emotional reaction \citep{hochschild2019managed}.  The stress associated with this type of emotional labor is often cited as a key cause of burnout and attrition among customer service workers \citep{leePhilippinesHasBecome2015}.  Additionally, contact center work for US-based businesses is frequently outsourced to countries such as India and the Philippines, meaning agents often work difficult hours and may face cultural barriers or judgments when speaking with customers.  Moreover, increases in productivity may not improve the work environment, especially if workers feel pressure to work faster. 

%In this section, we examine the impact of generative AI on two particularly salient aspects of the workplace experience: how agents are treated by customers and how often they face requests to speak with a manager.  We also examine the impact of AI assistance on worker turnover as an overall indicator of worker satisfaction. 

\subsection{Customer Sentiment}

Qualitative studies suggest that working conditions for contact center agents can be unpleasant.  Customers often vent their frustrations to anonymous service agents and, in our data, we see regular instances of swearing, verbal abuse, and ``yelling'' (typing in all caps).  The stress associated with this type of emotional labor is often cited as a key cause of burnout and attrition among customer service workers \citep{leePhilippinesHasBecome2015}.  

%Moreover, increases in productivity may not improve the work environment, especially if workers feel pressure to work faster. 
% Additionally, contact center work for US-based businesses is frequently outsourced to countries such as India and the Philippines, meaning agents often work difficult hours and may face cultural barriers or judgments when speaking with customers. 

Access to AI-assistance may impact how customers treat agents, but, in theory, the direction and magnitude of these impacts are ambiguous.  AI assistance may improve the tenor of conversations by helping agents set customer expectations or resolve their problems more quickly.  Alternatively, customers may become more frustrated if AI-suggested language feels ``corporate'' or insincere.  

To assess this, we capture the affective nature of both agent and customer text, using sentiment analysis \citep{mejovaSentimentAnalysisOverview}. For this analysis, we use SiEBERT, an LLM that is fine-tuned for sentiment analysis using a variety of datasets, including product reviews and tweets \citep{hartmann2023}.  Sentiment is measured on a scale from $-1$ to $1$, where $-1$ indicates negative sentiment and 1 indicates positive.  In a given conversation, we compute separate sentiment scores for both agent and customer text.  We then aggregate these chat-level variables into a measure of average agent sentiment and average customer sentiment for each agent-year-month.

Panels A and B of Figure \ref{fig:experienceofwork} consider how sentiment scores respond following the rollout of AI assistance.  In Panel A, we see an immediate and persistent improvement in customer sentiment.  This effect is large: According to Column 1 of Table \ref{tab:sentiment}, access to AI improves the mean customer sentiments (averaged over an agent-month) by 0.18 points, equivalent to half of a standard deviation.  In Panel B, we see no detectable effect for agent sentiment, which is already very high at baseline.  Column 2 of Table \ref{tab:sentiment} indicates agent sentiments increase by only 0.02 points or about 1 percent of a standard deviation. 

Focusing on customer sentiment, Panels C and D of Appendix Figure \ref{afig:expwork} examine whether access to AI has different impacts across agents.  We find access to AI assistance significantly improves how customers treat agents of all skill and experience levels, with the largest effects for agents in the lower to lower-middle range of both the skill and tenure distributions.  Consistent with our productivity results, the highest-performing and most-experienced agents see the smallest benefits of AI access.  These results suggest that AI recommendations, which were explicitly designed to prioritize more empathetic responses, may improve agents' demonstrated social skills and have a positive emotional impact on customers.  

\subsection{Customer Confidence and Managerial Escalation}

Changes in individual productivity may have broader implications for organizational workflows \citep{athey_allocation_1994, atheystern_1998, garicano_hierarchies_2000}.  
In most customer service settings, front-line agents attempt to resolve customer problems but can seek the help of supervisors when they are unsure how to proceed.  Customers, knowing this, will sometimes attempt to escalate a conversation by asking to speak to a manager. This type of request generally occurs when frustrated customers feel that the current agent is not adequately addressing their problem.

In Panel C of Figure \ref{fig:experienceofwork}, we consider the impact of AI assistance on the frequency of chat escalation.  The outcome variable we focus on is the share of an agent's chats in which a customer requests to speak to a manager or supervisor, aggregated to the year-month level.  We focus on requests for escalation rather than actual escalations both because we lack data on actual escalations and because requests are a better measure of customer confidence in an agent's competence or authority.  

Following the introduction of AI assistance, we see a gradual but substantial decline in requests for escalation.  Relative to a baseline rate of approximately 6 percent, these coefficients suggest that AI assistance generates an almost 25 percent decline in customer requests to speak to a manager.  In Panels E and F of Appendix Figure \ref{afig:expwork}, we consider how these patterns change depending on the skill and experience of the worker.  While these results are relatively noisy, our point estimates suggest that requests for escalation are disproportionately reduced for agents who were less skilled or less experienced at the time of AI adoption.   

\subsection{Attrition}

The adoption of generative AI tools can affect workers in various ways, including their productivity, the level of stress they experience, how customers perceive them, and their overall job satisfaction. While we cannot directly observe all these factors, we can analyze turnover patterns as a broad measure of how workers respond to AI implementation.

We compare attrition rates between AI-assisted agents and untreated agents with equal tenure. We drop observations for treated agents before treatment because they do not experience attrition by construction (they must survive to be treated in the future), and control for location and time fixed effects.
 
Consistent with our findings so far, Panel A of Appendix Figure \ref{afig:attrition} shows that access to AI assistance is associated with the strongest reductions in attrition among newer agents, those with less than 6 months of experience.  The magnitude of this coefficient, around 10 percentage points, translates into a 40 percent decrease relative to a baseline attrition rate in this group of 25 percent.  In Panel B, we examine attrition by worker skill.  Here, we find a significant decrease in attrition for all skill groups, although without a clear gradient.

These results should be taken with more caution relative to our main results because attrition occurs once per worker and therefore we are unable to include agent fixed effects. Our results may overstate the impact of AI access on attrition if, for example, the firm is more likely to give AI access to agents deemed more likely to stay.  


\section{Conclusion}

% 0.  BROAD INTRO AND SUMMARY OF FINDINGS.
% 1.  CAVEAT ONE FIRM ONE TIME
% ONE OCCUPATION THAT IS LOWER SKILL.  

% 2.  TALK ABOUT HOW THIS IS NOT JUST RAINBOWS FOR WORKERS
%  --- CONVERSATIONS WHERE THEY SAY THEYRE HIRING FEWER PEOPLE
%  --- NO EVIDENCE OF DE-SKILLING YET.  ORG INERTIA.  

% 3.  HARMS INCENTIVES FOR WORKER INNOVATION AND FUTURE MODEL TRAINING.  GROUPED WITH PAY COMMENTS.  
%  --- ALSO USE OF HIGHER SKILL WORKERS TO VALIDATE LLM RESPONSES---NEW JOB TYPES


Advances in AI technologies open up a broad set of economic possibilities.  Our paper provides early empirical evidence on the effects of a generative AI tool in a real-world workplace. In our setting, we find access to AI-generated recommendations increases overall worker productivity by 15 percent, with even larger impacts for lower-skill and novice agents. These productivity gains in part reflect durable worker learning rather than rote reliance on AI suggestions.  Furthermore, AI assistance appears to improve worker on-the-job experiences, such as by improving customer sentiment and confidence, and is associated with reductions in turnover.

Our analysis is subject to some caveats and raises many unanswered questions.

First, we note again that our findings apply for a particular AI tool, used in a single firm, within a single occupation, and should not be generalized across all occupations and AI systems.  For example, our setting has a relatively stable product and set of technical support questions. In areas where the product or environment is changing rapidly, the relative value of AI recommendations may be different. For instance, AI may be better able to synthesize changing best practices, or could impede learning by promoting outdated practices observed in historical training data.  Indeed, recent work  by \citet{Perry2022DoUW} and \citet{otisUnevenImpactGenerative2023} have found instances in which AI adoption has limited or even negative effects.  
%Our findings highlight the potential productivity benefits of AI, but underscore the need for broader analysis across various worker types, industries, and AI systems.

Second, we report partial equilibrium short- to medium-run impacts of AI deployment.  While we do not have access to compensation data, the managers we spoke to believed that workers may have received higher performance pay as a result of AI assistance, since these bonuses were typically tied to targets related to average handle time and resolution rates. They caution, however, that potential gains in bonus pay may not be long-lived because it is common practice to adjust performance targets if too many agents were hitting the goals.  As a result, workers may eventually be subject to a ratchet effect if AI assistance leads performance targets to be readjusted upward.  

More generally, we are not able to observe longer-run equilibrium responses.  In principle, the increased productivity we observe could lead to either lower or higher demand for customer service agents. If customer demand for assistance is inelastic, then the productivity gains we document will likely translate into less demand for human labor. A back-of-the-envelope calculation suggests the firm could field the same number of customer support issues with 12 percent fewer worker-hours.  Conversely, individuals may currently avoid contacting customer service because of the long wait times and low-quality service.  AI assistance that improves this experience may boost consumer demand for product support, resulting in increased labor demand \citep{berg2018, korinekantonHowInnovationAffects2023}.  In addition, the use of AI could create new jobs for customer service agents, such as testing and training AI models \citep{autornewwork2022}. One manager we spoke with reports that high-skill workers in some contact centers are already being tasked with reviewing AI suggestions and providing better alternatives.  Other work shows that even low levels of AI adoption can impact market equilibrium prices and quantities, highlighting the need for more work on the equilibrium effects of AI on the labor market \citep{raymondjmp}.

%Furthermore, the eventual impacts of AI tools on workers, wages, and aggregate productivity effects will depend on context-specific factors and the development of complementary organizational practices and educational policies \citep{brynjcurve}.  For example, \citet{humlumAdoptionChatGPT2024} show workers' willingness to use AI depends critically on access to training and on regulations such as those related to data confidentiality.  

Finally, our findings also raise questions about the nature of worker productivity.  Traditionally, a support agent's productivity refers to their ability to help the customers.  Yet, in a setting where customer service conversations are fed into training datasets, a worker's productivity also includes the AI training data they produce.  Top performers, in particular, contribute many of the examples used to train the AI system we study.  This increases their value to the firm. At the same time, our results suggest that access to AI suggestions may lead them to put less effort into coming up with new solutions.  Going forward, compensation policies that provide incentives for people to contribute to model training could be important. Given the early stage of generative AI, these and other questions deserve further scrutiny.  
\vspace{10pt}

\noindent\textsc{Stanford University and National Bureau of Economic Research}\\
\textsc{Massachusetts Institute of Technology and National Bureau of Economic Research}\\
\textsc{Massachusetts Institute of Technology}
% \begin{tabular}{ccccccc}
% \large{Erik Brynjolfsson} & \mbox{ } & \mbox{ } & \large{Danielle Li} & \mbox{ } & \mbox{ } & \large{Lindsey Raymond} \\
% \large{Stanford \& NBER} & \mbox{ } & \mbox{ } & \large{MIT \& NBER} & \mbox{ } & \mbox{ } & \large{MIT} \\
% \end{tabular}
\section*{Supplementary Material}
The Appendix for this article can be found at \textit{The Quarterly Journal of Economics} online. 
\clearpage
\begin{spacing}{0.8}
%\bibliographystyle{aea}
% \bibliography{AugI_Bibliography.bib}
\bibliography{cleaned_bibliography}
\end{spacing}

%%%%%%%%%%%%%%%%%%%%%%%%%%%%%%%%%%%%%%%%%%%%%%%%%%%%%%%
%%%%%%%%%%%%%%%%%%%%%%%%%%%%%%%%%%%%%%%%%%%%%%%%%%%%%%%
%TABLES AND FIGURES IN ANOTHER FILE, INPUTTED HERE.
%%%%%%%%%%%%%%%%%%%%%%%%%%%%%%%%%%%%%%%%%%%%%%%%%%%%%%%
%%%%%%%%%%%%%%%%%%%%%%%%%%%%%%%%%%%%%%%%%%%%%%%%%%%%%%%

\input{tablesandfigs.tex}

\end{document}